\section{Responsible deployment}
\label{sec:responsible}

\subsection{Defensive scope}

\cyberlite{} is framed for defensive security: detection engineering, secure
coding, vulnerability remediation, hardening, and authorized education. It is not
an autonomous exploitation system. The repository does not provide tools for
unauthorized access, and examples emphasize authorization, evidence handling, and
verification.

\subsection{No execution}

Generated commands and code are returned as text. The inference, autocomplete,
evaluation, and server modules contain no shell or tool-execution path.
Autocomplete responses explicitly report \texttt{executed: false} and
\texttt{requires\_review: true}. Evaluation checks Python syntax with
\texttt{ast.parse} but never runs generated code. Integrators must add a sandbox,
authorization checks, and human review if they choose to execute anything.

\subsection{Human oversight}

Security advice is consequential. A fluent answer can be wrong, outdated, or
unsafe. Deployments should:

\begin{itemize}[leftmargin=1.4em]
  \item present outputs as suggestions, not decisions;
  \item require verification against authoritative sources and tested
  configurations;
  \item log model/version/prompt for auditability without logging secrets;
  \item provide an escalation path to a human expert;
  \item restrict use to authorized environments.
\end{itemize}

\subsection{Privacy and data handling}

Local inference reduces exposure to third-party providers but does not eliminate
risk. Prompts, logs, feedback files, and evaluation outputs can contain sensitive
data. The feedback endpoint stores hashes rather than completion text by default,
but operators must configure retention and redaction. Do not send production
credentials, personal data, or customer telemetry to a model without policy
approval.

\subsection{Dual use}

Cybersecurity knowledge is dual-use. The training corpus may contain descriptions
of attacks, vulnerabilities, and malware. Defensive framing reduces but does not
eliminate the risk of harmful recombination. The project supports evaluation of
refusal and over-refusal, and it encourages red-teaming before release. Model
cards and dataset cards should state intended and out-of-scope use
\citep{mitchell2019modelcards}.

\subsection{Licensing and attribution}

Code is Apache-2.0. Model weights and datasets retain their own licenses. ATT\&CK
and CWE require attribution under their terms; OWASP content may be share-alike;
SecEval is evaluation-only and non-commercial; CyberSecEval is obtained under
upstream terms. The project does not redistribute restricted data. Users are
responsible for compliance when publishing derived datasets or models.

\subsection{Disclosure}

Security vulnerabilities in the repository should be reported privately according
to \texttt{SECURITY.md}. Vulnerabilities in upstream models, libraries, or
datasets should be reported to their maintainers. The project will document
known limitations rather than claiming safety from a smoke test.

\subsection{Deployment checklist}

Before exposing a \cyberlite{} service:

\begin{enumerate}[leftmargin=1.4em]
  \item bind to loopback or authenticated, encrypted access;
  \item set an API key and rate limits if other users can reach it;
  \item redact or disable prompt-body logging;
  \item never connect output to a shell, CI runner, or production system;
  \item monitor memory, latency, and abuse;
  \item publish the model card, dataset card, evaluation protocol, and
  limitations;
  \item re-evaluate after model, data, or dependency changes.
\end{enumerate}
